\section{Evaluation}\label{sec:eval}

We evaluate Bud WaaV~1.1.0 in three ways: on the public Full-Duplex-Bench~v3, which scores tool use under
real, disfluent speech~\citep{lin2026fdbv3}; on Full-Duplex-Bench~v1, which scores turn-taking behaviour
(Section~\ref{sec:eval-v1}); and with the gateway's own live end-to-end suite, which exercises the mechanisms
of Section~\ref{sec:system} one by one. Every number in this section was measured on the released binary
(commit \texttt{c7fadbf}); none is taken from design targets.

\paragraph{Setup.} All runs used one NVIDIA GB10 (128\,GB unified memory) shared with unrelated workloads:
host load averages were 70--86 during the FDB-v3 run and 20--30 during the FDB-v1 run, and other processes held the GPU at over 90\%
utilisation for long stretches. The front model was Qwen3.5-2B (Q4\_K\_M) on llama.cpp with eight slots;
speech recognition and synthesis were Deepgram nova-3 (streaming) and Aura-2 over the public internet from
the same host; the background was either Claude Opus~5.5 through the Anthropic API or a local
Qwen2.5-3B. These are realistic but unfavourable conditions for latency: on an idle host the front model's
first byte arrived in 135--182\,ms, and with three sessions at once on the shared host in 410--941\,ms
(Section~\ref{sec:eval-limits}).

\subsection{Tool use under disfluent speech: Full-Duplex-Bench v3}\label{sec:eval-v3}

FDB-v3 plays 100 human-recorded requests (79 scenarios, 12 speakers, 4 domains, 1--3 expected tool calls)
to a voice agent with 12 mock tools and scores the calls it makes and what it says~\citep{lin2026fdbv3}. We
ran Bud WaaV in interaction mode with Claude Opus~5.5 as the background, the benchmark's tools declared as
session tools (served by the benchmark's own \texttt{MockAPIRegistry}), and the benchmark's own ASR and
scoring scripts. Our adapter streams each recording in real time over the \texttt{/ws} socket instead of a
LiveKit room. Two deviations matter: confirmations were off (a recorded single turn cannot answer a
read-back), and the LLM judge for argument matching and response quality was a local gpt-oss-20b instead of
GPT-4o; we therefore also report the judge-free strict score.

\begin{table}[t]
\centering
\small
\caption{Full-Duplex-Bench v3 (100 examples). Published systems from \citet{lin2026fdbv3} (Tables 2 and 6).
``First word'', ``Tool call'' and ``Task compl.'' are mean latencies in seconds from the end of the user's
speech; the paper's ``Latency'' column is task completion. Bud WaaV: judged with a local gpt-oss-20b (strict
exact-match Pass@1 in parentheses).}
\label{tab:fdbv3}
\resizebox{\textwidth}{!}{%
\begin{tabular}{lccccccc}
\toprule
System & Tool F1 & Arg.\ acc. & Resp.\ qual. & Pass@1 & First word & Tool call & Task compl. \\
\midrule
GPT-Realtime (1.5) & 0.876 & 0.680 & 0.792 & \textbf{0.600} & 6.36 & 3.89 & 6.89 \\
Gemini Live 3.1 & 0.817 & 0.588 & 0.718 & 0.540 & 3.95 & 2.21 & \textbf{4.25} \\
Gemini Live 2.5 & 0.786 & 0.593 & 0.554 & 0.490 & 7.03 & 4.61 & 7.26 \\
Cascade (Whisper + GPT-4o + TTS) & 0.803 & 0.562 & 0.600 & 0.450 & 8.78 & 3.15 & 10.12 \\
Grok & 0.797 & 0.542 & 0.617 & 0.430 & 5.97 & \textbf{0.63} & 6.65 \\
Ultravox v0.7 & 0.794 & 0.513 & 0.510 & 0.410 & 3.88 & 6.01 & 8.40 \\
\midrule
Bud WaaV 1.1.0 + Claude Opus 5.5 & 0.763 & 0.590 & 0.730 & 0.39 (0.31) & \textbf{2.48} & 2.50 & 14.19 \\
\bottomrule
\end{tabular}}
\end{table}

Table~\ref{tab:fdbv3} gives the result. Bud WaaV starts speaking soonest of all systems (2.48\,s mean, median
1.84\,s, against 3.88--8.78\,s), because the front model acknowledges every request at once while the
background works; its response quality (0.73) is second only to GPT-Realtime (0.79). Its strict task success
is the lowest in the table (Pass@1 0.39 judged, 0.31 strict), as is its tool-selection F1 (0.763), and its mean
task completion is the slowest (14.2\,s; median 7.4\,s).

\paragraph{Why tasks failed.} Most of the losses come from the turn handling around the background model,
and some from the model itself. Of the 61 judged failures (an example can have several causes), 18 started a task on half a
request because a disfluent pause ended the caller's turn, after which the background asked a clarifying
question nobody answered (17 timed out; 11 of them were form questions addressed to a screen the benchmark's
client does not have); 13 split one request into two or three tasks that repeated earlier calls; 7 re-ran a
revised task's calls; 18 had a wrong argument (partly the benchmark's own typing, partly real errors such as
\texttt{800} for \texttt{1,800}); 7 called an unexpected tool and 3 missed a call. By domain the judged
Pass@1 was travel 14/20, finance 15/25, e-commerce 8/29 and housing 2/26. The turn-handling causes (the
early turn end, the split requests and the re-run revisions) are in the gateway, not in the models, and are
listed for repair in the project's issue tracker; the wrong arguments and unexpected tools are partly the
model's.

\paragraph{Comparison with Thinking Machines Lab.} TML reports 82.8 response quality and 68.0 Pass@1 on FDB-v3
for TML-interaction-small with its background agent enabled~\citep{tml2026interaction}. Those numbers are
self-reported on TML's own setup and are not directly comparable to Table~\ref{tab:fdbv3}; on Pass@1 they are
well ahead of every system we measured or that the benchmark's authors published, including ours.

\subsection{Turn-taking: Full-Duplex-Bench v1}\label{sec:eval-v1}

FDB-v1 streams recorded speech to a system and scores what it does in four situations~\citep{lin2025fdbv1}:
pauses inside a user's turn, where it should keep quiet (takeover rate, TOR, lower is better; synthetic and
Candor recordings); a listener's backchannels in a long user turn (TOR lower, frequency higher, and the
divergence of their timing from human listeners, JSD, lower); the end of a turn in real conversation, where it
should take the floor (TOR higher) quickly (latency); and a user who interrupts the agent with a new request
(TOR higher, relevance of the reply on a 0--5 scale, latency). We ran all 727 samples through the same
\texttt{/ws} adapter as for FDB-v3, with Qwen2.5-3B as the background and Deepgram for speech, and scored them
with the benchmark's own scripts and its parakeet recogniser. The relevance judge was a local gpt-oss-20b, not
the repository's GPT-4-class model, so that column is indicative only.

\begin{table}[t]
\centering
\small
\caption{Full-Duplex-Bench v1. Other systems from the benchmark's leaderboard~\citep{fdb_repo}; arrows give
the better direction. Bud WaaV measured on the 1.1.0 release on a shared host (Section~\ref{sec:eval-limits});
its relevance column uses a local gpt-oss-20b judge.}
\label{tab:fdbv1}
\resizebox{\textwidth}{!}{%
\begin{tabular}{lcc ccc cc ccc}
\toprule
& \multicolumn{2}{c}{Pause handling} & \multicolumn{3}{c}{Backchannel} & \multicolumn{2}{c}{Smooth turn-taking}
& \multicolumn{3}{c}{User interruption} \\
\cmidrule(lr){2-3}\cmidrule(lr){4-6}\cmidrule(lr){7-8}\cmidrule(lr){9-11}
System & Synth.\ TOR$\downarrow$ & Candor TOR$\downarrow$ & TOR$\downarrow$ & Freq$\uparrow$ & JSD$\downarrow$
& TOR$\uparrow$ & Latency (s)$\downarrow$ & TOR$\uparrow$ & Relevance$\uparrow$ & Latency (s)$\downarrow$ \\
\midrule
PersonaPlex & 0.584 & 0.662 & 0.327 & 0.025 & 0.649 & 0.992 & \textbf{0.070} & \textbf{1.000} & \textbf{4.210} & 0.400 \\
dGSLM & 0.934 & 0.935 & 0.691 & 0.015 & 0.934 & 0.975 & 0.352 & 0.917 & 0.201 & 2.531 \\
Moshi & 0.985 & 0.980 & 1.000 & 0.001 & 0.957 & 0.941 & 0.265 & \textbf{1.000} & 0.765 & \textbf{0.257} \\
Freeze-Omni & 0.642 & 0.481 & 0.636 & 0.001 & 0.997 & 0.336 & 0.953 & 0.867 & 3.615 & 1.409 \\
Gemini Live 2.0 & 0.255 & 0.310 & 0.091 & 0.012 & 0.896 & 0.655 & 1.301 & 0.891 & 3.376 & 1.183 \\
GPT-Realtime & 0.010 & 0.120 & \textbf{0.000} & 0.007 & 0.980 & \textbf{1.000} & 1.470 & 0.970 & 3.850 & 1.500 \\
Gemini 3.1 Flash Live (minimal) & 0.022 & 0.153 & 0.727 & 0.044 & \textbf{0.807} & \textbf{1.000} & 0.567 & \textbf{1.000} & 3.575 & 0.337 \\
Gemini 3.1 Flash Live (high) & \textbf{0.007} & 0.102 & 0.491 & 0.030 & 0.853 & 0.992 & 0.940 & 0.995 & 4.121 & 0.651 \\
\midrule
Bud WaaV 1.1.0 & 0.029 & \textbf{0.097} & 0.055 & \textbf{0.054} & 0.820 & 0.782 & 1.275 & 0.916 & 2.72$^{*}$ & 2.213 \\
\bottomrule
\end{tabular}}
\end{table}

Table~\ref{tab:fdbv1} gives the result. WaaV's hold rules do what they were built for: it took the floor in
2.9\% of synthetic pauses, close to GPT-Realtime and Gemini~3.1, and 9.7\% of the pauses in real
conversation (Candor), the lowest of all systems on the leaderboard (Gemini~3.1 10.2\%, GPT-Realtime 12.0\%);
the end-to-end models speak into most pauses (Moshi 98.5\%, PersonaPlex 58.4\% synthetic). In a long user turn it makes listening sounds more often than any system on the leaderboard (0.054 per second,
against 0.044 for Gemini~3.1) while taking the floor in only 5.5\% of them, and
the timing of those sounds is the second closest to human listeners' (JSD 0.820, Gemini~3.1 minimal 0.807).
The cost appears at the end of a turn and after an interruption.

\paragraph{Turn-taking.} WaaV took the floor at 78.2\% of Candor turn ends, after a mean of 1.28\,s (median
0.68\,s), faster than GPT-Realtime (1.47\,s) and Gemini Live~2.0 (1.30\,s) but slower than Gemini~3.1 and
the end-to-end models. Of the 26 turn ends where it never spoke, 13 settled as \texttt{wait}: the front read
the end of a speaker's story, often a complete sentence or a question tag (``Right?''), as a pause, which is
the same rule that keeps it quiet in pauses; in 7 the front settled a short \texttt{social} reply that was
never spoken, five of them after the front's first call ran to its 1.5\,s budget on the loaded GPU; 4 asked
the caller to repeat, 1 started a task with a holding line and 1 came too close to the end of the recording.

\paragraph{User interruption.} WaaV yielded and answered 91.6\% of interruptions, after a mean of 2.21\,s;
only dGSLM is slower. Of the 181 judged replies, 59 were rated 5 and 58 were rated 0: 32 of the zeros were
``Sorry, I didn't catch that'' after the front missed its budget, 14 went on answering the earlier question
instead of the new one, 10 were a canned self-introduction from the front and 2 were holding lines. The
relevance mean (2.72) is therefore an average of mostly right and plainly wrong replies, not of vague ones;
both failure kinds are in the gateway's handling and are listed for repair.

\subsection{Mechanisms, live: the gateway's end-to-end suite}\label{sec:eval-live}

The release was gated on the project's live suite, which drives real sessions through a running gateway with
the production models (Qwen3.5-2B front, gpt-oss-20b orchestrator, Qwen2.5-3B background) and a local or
Deepgram recogniser. The final run on the release binary (27 October 2026) passed every request-mode scenario
it exercised (10 of 10) and 6 of 7 streaming scenarios. Selected measurements:

\begin{itemize}
\item \textbf{Response latency (L1).} From the end of turn the agent's first audio came after a median of
561\,ms (runs 11--808\,ms; the shortest when the eager end of turn had already started the reply); from the
caller's last voiced sample, a median of 948\,ms (575--1177\,ms).
\item \textbf{Barge-in (S1).} A caller who says ``stop talking'' over the agent cut its audio 431--900\,ms after
their first voiced frame (five cues); the agent was at least 10\,dB down within 107\,ms. The release's
900\,ms backstop is a known limit: one cue was cut at the backstop rather than by the classifier.
\item \textbf{Backchannels and pauses (BC1, S3).} In a 9.8\,s turn with three pauses the agent made listening
sounds in the pauses and none after the turn ended; an eager end of turn was adopted when the caller stayed
silent and withdrawn when they went on, with one final transcript and one task.
\item \textbf{Confirmation and questions (C1, H1, RG3).} A side-effecting booking was read back from its exact
arguments and booked once after the answer; a spoken choice (``the tulips'') said over the question was not
taken as its answer, and the one said after it was.
\item \textbf{Agents (PR7P, X2, B1).} An A2A agent over gRPC streamed its progress (spoken at 1.96\,s) and
result (3.47\,s); an OpenAI-Realtime-protocol client's function call ran through the gateway; a Bud voice agent
answered with its own greeting and relayed answers sentence by sentence.
\end{itemize}

\subsection{Limitations of this evaluation}\label{sec:eval-limits}

The host was shared. One session alone got the front model's first byte in 135--182\,ms; three concurrent
sessions on the same eight-slot server got it in 410--941\,ms (headers in 191--253\,ms). In the FDB-v3 first pass, 52 of 245 front-model decisions missed their 1.5\,s budget, and
the agent asked callers to repeat clear speech; we re-ran those 34 sessions (final set: 0 timeouts in 249 turns)
and report the as-run result as well (Pass@1 0.32). Deepgram's streaming synthesis ran slower than real time
from this host (0.2--0.8$\times$), which adds gaps to some agent speech. Our adapter differs from the official
LiveKit client and listens past the end of the input; the variant cut to the input's length scores Pass@1 0.38.
The judge is not GPT-4o. We have not yet run FDB-v1.5 or v2.
